\begin{kbvderiv}
\vspace{4pt plus 3pt minus 1pt}
\noindent{\renewcommand{\arraystretch}{1.24}\setlength{\tabcolsep}{5pt}
\begin{tabularx}{\textwidth}{X X}
ಮಾಡು + ಅದು = ಮಾಡದು & ಮಾಡು + ಅವು = ಮಾಡವು \\
ಮಾಡು + ಎ = ಮಾಡೆ & ಮಾಡು + ಅರಿ = ಮಾಡಿರಿ \\
ಮಾಡು + ಎನು = ಮಾಡೆನು & ಮಾಡು + ಎವು = ಮಾಡೆವು \\
\end{tabularx}}
\vspace{4pt plus 3pt minus 1pt}
\end{kbvderiv}

\indent ಇವು “ಮಾಡು” ಧಾತುವಿನ ಮೇಲೆ ಹಲವು ಪ್ರತ್ಯಯಗಳು ಸೇರಿ “ಕ್ರಿಯೆಯು ನಡೆಯದು” ಎಂಬ ಅರ್ಥ ಕೊಡುವ ಕ್ರಿಯಾಪದಗಳಾಗಿರುವುವು. ಹೀಗೆ, ಕ್ರಿಯೆಯು ನಡೆಯದು ಎಂಬರ್ಥ ಕೊಡುವ ಕ್ರಿಯಾಪದಗಳಿಗೆ ನಿಷೇಧಾರ್ಥಕ ಕ್ರಿಯಾಪದಗಳೆಂದು ಹೆಸರು.\par

\indent : ಧಾತುಗಳ ಮೇಲೆ ಏಕವಚನದಲ್ಲಿ ಅನು, ಅಳು, ಅದು, ಎ, ಎನು ಎಂಬ ಆಖ್ಯಾತ ಪ್ರತ್ಯಯಗಳೂ ಬಹುವಚನದಲ್ಲಿ ಅರು, ಅವು, ಅರಿ, ಎವು ಎಂಬ ಆಖ್ಯಾತ ಪ್ರತ್ಯಯಗಳೂ ಸೇರಿ ನಿಷೇಧಾರ್ಥಕ ಕ್ರಿಯಾಪದಗಳಾಗುವುವು.\par

\vspace{2.5pt plus 3pt minus 0.5pt}\noindent\hspace*{1.2em}\kbvnum{(6)}\hspace{0.55em}ಸಂಭಾವನಾರ್ಥಕ ಕ್ರಿಯಾಪದಗಳು\par

\noindent\hspace*{1.2em}\kbvlabel{ಉದಾ}\hspace{0.6em}“ಕೊಡು” ಧಾತು\par

\vspace{5pt plus 4pt minus 1pt}
\begin{center}
{\renewcommand{\arraystretch}{1.24}\setlength{\tabcolsep}{10pt}\rowcolors{2}{kbvrow}{white}
\begin{tabularx}{0.88\textwidth}{X X}
{\kbvsecfont\color{kbvaccentdark}ಏಕವಚನ} & {\kbvsecfont\color{kbvaccentdark}ಬಹುವಚನ} \\[1pt]
\arrayrulecolor{kbvline}\specialrule{0.5pt}{0pt}{2pt}\arrayrulecolor{black}
ಕೊಡು + ಆನು = ಕೊಟ್ಟಾನು & ಕೊಡು + ಆರು = ಕೊಟ್ಟಾರು \\
ಕೊಡು + ಆಳು = ಕೊಟ್ಟಾಳು & ಕೊಡು + ಆರು = ಕೊಟ್ಟಾರು \\
ಕೊಡು + ಈತು = ಕೊಟ್ಟೀತು & ಕೊಡು + ಆವು = ಕೊಟ್ಟಾವು \\
ಕೊಡು + ಈಯೇ = ಕೊಟ್ಟೀಯೆ & ಕೊಡು + ಈರಿ = ಕೊಟ್ಟೀರಿ \\
ಕೊಡು + ಏನು = ಕೊಟ್ಟೇನು & ಕೊಡು + ಏವು = ಕೊಟ್ಟೇವು \\
\end{tabularx}}
\end{center}
\vspace{5pt plus 4pt minus 1pt}

\indent {[}ಗಮನಿಸಿ : “ಕೊಡು” ಧಾತುವಿನ ಮೇಲೆ ಸಂಭಾವನಾರ್ಥಕ ಆಖ್ಯಾತ ಪ್ರತ್ಯಯಗಳು ಸೇರಿದಾಗ ಧಾತುವಿನ ಕೊನೆಯ “ಡು” ಕಾರಕ್ಕೆ “ಟ್ಟ” ಕಾರ ಆದೇಶವಾಗಿ ಬರುವುದು.{]}\par

\indent ಇಲ್ಲಿ “ಕೊಡು” ಧಾತುವಿನ ಮೇಲೆ ಹಲವು ಪ್ರತ್ಯಯಗಳು ಸೇರಿ “ಕ್ರಿಯೆಯು ಸಂಭವಿಸುವಲ್ಲಿ ಸಂಶಯತೋರುವ” ಕ್ರಿಯಾಪದಗಳಾಗಿರುವುವು.\par
